\section{谱}
\label{sec:ospec-spectrum}

在$n$维复空间中，$\bm A-\lambda\bm I$不可逆当且仅当它不是单射，即存在非零向量$\bm v$使$\bm A\bm v=\lambda\bm v$。有限维中单射、满射与可逆三者等价，所以“使$\bm A-\lambda\bm I$不可逆的$\lambda$”恰好就是特征值；特征值与对称矩阵的正交对角化可参考\BookChapterRef{chap:eigen-spectral-decomposition}{卷二第4章“特征值与谱分解”}。无限维中三者不再等价，算子可以是单射而不满射，因此必须直接从可逆性出发。

本节设$X\neq\{0\}$是复Banach空间，$T\in\mathcal B(X)$。要求$T$回到同一个空间，是因为$T-\lambda I$只对这样的算子有意义。实空间上的算子可以先复化：令$X_{\mathbb C}=\{x+\mathrm iy:x,y\in X\}$，按复数规则定义数乘，令$T_{\mathbb C}(x+\mathrm iy)=Tx+\mathrm iTy$；Hilbert空间中取$\lVert x+\mathrm iy\rVert^2=\lVert x\rVert^2+\lVert y\rVert^2$，它来自内积$\langle x+\mathrm iy,u+\mathrm iv\rangle=\langle x,u\rangle+\langle y,v\rangle+\mathrm i(\langle x,v\rangle-\langle y,u\rangle)$。实空间算子的谱就指$T_{\mathbb C}$的谱。

复数域是必要的。平面上旋转$90^\circ$的实矩阵$\left[\begin{smallmatrix}0&-1\\1&0\end{smallmatrix}\right]$没有实特征值，复化后才有特征值$\pm\mathrm i$。紧算子的特殊性质留到下一节，本节只讨论一般有界算子。

\subsection{谱与预解集}
\label{sec:ospec-resolvent}

参数$\lambda$“好”的含义是：对每个右端$g\in X$，方程$(T-\lambda I)u=g$都有解（存在），解只有一个（唯一），并且$g$的小误差只引起$u$的小误差（稳定）。三条合起来，就是$T-\lambda I$有定义在整个$X$上的有界逆。

\begin{definition}[谱与预解集]
\label{def:ospec-spectrum}
若$T-\lambda I$是$X$到$X$的双射，且逆映射$(T-\lambda I)^{-1}$有界，则称$\lambda\in\mathbb C$属于$T$的\term{预解集}（resolvent set），记为$\rho(T)$，并称$R_\lambda=(T-\lambda I)^{-1}$为\term{预解算子}（resolvent operator）。预解集的补集$\sigma(T)=\mathbb C\setminus\rho(T)$称为$T$的\term{谱}（spectrum）。若存在非零$u\in X$使$Tu=\lambda u$，则称$\lambda$为$T$的\term{特征值}（eigenvalue），$u$为相应的\term{特征向量}（eigenvector），函数空间中也称特征函数。
\end{definition}

由第~\ref{chap:linear-operators}章的有界逆定理，Banach空间之间的有界线性双射自动有有界逆，所以定义中“逆有界”一项不额外增加限制；写出它，是为了强调预解参数要求在整个$X$上稳定地反解。定义也没有要求事先找到特征向量：$\lambda$进入谱，可以因为$T-\lambda I$有非零零空间（$\lambda$是特征值），也可以因为它是单射却不满射，此时方程对某些右端无解。

\paragraph{谱的有界性与闭性}
证明需要一个反解工具：若$\lVert A\rVert<1$，则$I-A$可逆。直观上，$(I-A)^{-1}$应当是几何级数$I+A+A^2+\cdots$；算子空间完备，这个级数确实收敛。

\begin{lemma}[小扰动的可逆性]
\label{lem:ospec-neumann}
设$A\in\mathcal B(X)$且$\lVert A\rVert<1$。则$I-A$有有界逆，$(I-A)^{-1}=\sum_{j\geq0}A^j$按算子范数收敛，且$\lVert(I-A)^{-1}\rVert\leq1/(1-\lVert A\rVert)$。
\end{lemma}

\begin{proof}
由复合的范数估计，$\lVert A^j\rVert\leq\lVert A\rVert^j$，所以$\sum_j\lVert A^j\rVert$是收敛的几何级数。$\mathcal B(X)$完备（第~\ref{chap:linear-operators}章），部分和$S_n=\sum_{j\leq n}A^j$按算子范数收敛到某个$S$，且$\lVert S\rVert\leq1/(1-\lVert A\rVert)$。直接相乘得$(I-A)S_n=S_n(I-A)=I-A^{n+1}$，令$n\to\infty$，由$\lVert A^{n+1}\rVert\to0$得$(I-A)S=S(I-A)=I$。
\end{proof}

这个级数及其截断误差在第~\ref{chap:linear-operator-equations}章作为求解第二类方程的方法系统讨论，这里只用它来确定谱的位置。

\begin{theorem}[谱是有界闭集]
\label{thm:ospec-spectrum-compact}
$\sigma(T)$是$\mathbb C$中的有界闭集，且$\sigma(T)\subseteq\{\lambda:\lvert\lambda\rvert\leq\lVert T\rVert\}$。对$\lambda\in\rho(T)$，有$\operatorname{dist}(\lambda,\sigma(T))\geq1/\lVert R_\lambda\rVert$。
\end{theorem}

\begin{proof}
有界性：若$\lvert\lambda\rvert>\lVert T\rVert$，写$T-\lambda I=-\lambda(I-T/\lambda)$，其中$\lVert T/\lambda\rVert<1$，由引理~\ref{lem:ospec-neumann}可逆，故$\lambda\in\rho(T)$。

闭性：设$\lambda_0\in\rho(T)$。对任意$\lambda$，
\[
T-\lambda I=(T-\lambda_0I)-(\lambda-\lambda_0)I=(T-\lambda_0I)\bigl(I-(\lambda-\lambda_0)R_{\lambda_0}\bigr).
\]
若$\lvert\lambda-\lambda_0\rvert<1/\lVert R_{\lambda_0}\rVert$，右端第二个因子由引理可逆，两个可逆有界算子的复合仍可逆，所以$\lambda\in\rho(T)$。这说明$\rho(T)$开，$\sigma(T)$闭，并且以$\lambda_0$为心、半径$1/\lVert R_{\lambda_0}\rVert$的开圆盘不含谱点。
\end{proof}

定理最后一式说明，参数越靠近谱，预解算子的范数至少要越大：反解$(T-\lambda I)u=g$的误差放大倍数在谱点附近没有上界。复Banach空间上有界算子的谱总是非空的，证明要用复变函数的Liouville定理考察预解算子关于$\lambda$的解析性，这里不给出证明；实空间不复化时，这一结论可以失败，旋转矩阵就是例子。

\subsection{谱点与特征值}
\label{sec:ospec-eigenvalues}

有限维中谱就是特征值的集合。下面的例子说明，无限维中谱可以完全不含特征值。

\paragraph{右移算子}
在$\ell^2$上令右移$S(x_1,x_2,\ldots)=(0,x_1,x_2,\ldots)$。它保持长度，$\lVert Sx\rVert=\lVert x\rVert$，所以$\lVert S\rVert=1$，并且$S$是单射。像中的序列首项必须为零，$e_1=(1,0,0,\ldots)$不在$\operatorname{Ran}S$中，所以$S$不满射，$S-0\cdot I$不可逆：$0\in\sigma(S)$。但$0$不是特征值，因为$Sx=0$只有零解。

$S$根本没有特征值。设$Sx=\lambda x$，逐坐标比较得$0=\lambda x_1$与$x_k=\lambda x_{k+1}$（$k\geq1$）。若$\lambda=0$，则$x_k=0$对一切$k$成立；若$\lambda\neq0$，由第一式$x_1=0$，再由第二式依次得$x_2=x_3=\cdots=0$。两种情形都只有零解。

谱却不小。由定理~\ref{thm:ospec-spectrum-compact}，$\sigma(S)$含于闭单位圆盘。对$\lvert\lambda\rvert<1$，第~\ref{chap:linear-operators}章给出$S^*$是左移，$(S-\lambda I)^*=S^*-\bar\lambda I$。序列$y=(1,\bar\lambda,\bar\lambda^2,\ldots)$平方可和，并满足$S^*y=(\bar\lambda,\bar\lambda^2,\ldots)=\bar\lambda y$，所以$\ker(S-\lambda I)^*\neq\{0\}$。由像与零空间的正交关系，$\operatorname{Ran}(S-\lambda I)$的闭包是$\ker(S-\lambda I)^*$的正交补，不是全空间，于是$S-\lambda I$不满射，$\lambda\in\sigma(S)$。谱是闭集，因而$\sigma(S)$是整个闭单位圆盘。

右移的谱充满一个圆盘，却没有一个特征向量。所以一般有界算子的谱不能理解为“特征值的集合”，后面基于特征向量的展开也不是所有有界算子都具备的表示。要得到特征展开，必须对算子再加条件，下一节的紧性是第一步。
